\section{Affine Connections}

\subsection{Levi-Civita Connections}
When it comes to connections, most of us probably think first of 
the Levi-Civita connection in Riemannian geometry. 
So let us first review its most basic definition, 
and then try to rewrite everything we are familiar with in the language of differential forms.
\begin{definition}
    A connection $\nabla$ on a mannifold $M$ is a map
    \begin{align*}
        \nabla: \Gamma(TM) \times \Gamma(TM) &\to \Gamma(TM) \\
        (X,Y) &\mapsto \nabla_X Y
    \end{align*}
    satisfying the following properties:
    \begin{align*}
        \nabla_{f_1X_1 + f_2X_2} Y =& f_1 \nabla_{X_1} Y + f_2 \nabla_{X_2} Y, \\
        \nabla_X (\lambda_1Y_1 + \lambda_2 Y_2) =& \lambda_1 \nabla_X Y_1 + \lambda_2 \nabla_X Y_2, \\
        \nabla_X (fY) =& \md f(X)\cdot Y + f \nabla_X Y. 
    \end{align*}
    where $X,X_1,X_2,Y,Y_1,Y_2 \in \Gamma(TM)$, $f,f_1,f_2 \in C^\infty(M)$, 
    and $\lambda_1,\lambda_2 \in \mathbb{R}$ are constants.
\end{definition}
\begin{definition}
    A connection $\nabla$ of a Riemannian manifold $(M,g)$ is called a Levi-Civita connection 
    if it is torsion-free and compatible with the metric $g$, i.e. 
    \begin{gather*}
        \nabla_X Y - \nabla_Y X = [X,Y], \\
        X(g(Y,Z)) = g(\nabla_X Y, Z) + g(Y, \nabla_X Z).
    \end{gather*}
\end{definition}
\begin{definition}
    The torsion of a connection $\nabla$ is defined as the tensor
    \begin{equation*}
        T(X,Y) := \nabla_X Y - \nabla_Y X - [X,Y].
    \end{equation*}
    The curvature of a connection $\nabla$ is defined as the tensor
    \begin{equation*}
        R(X,Y)Z := \nabla_X \nabla_Y Z - \nabla_Y \nabla_X Z - \nabla_{[X,Y]} Z.
    \end{equation*}
\end{definition}

\subsection{General Affine Connections}
\subsubsection{Horizontal Distribution}
\subsubsection{Connection 1-forms}
\subsubsection{Covariant Derivative}
\subsubsection{Curvature of Connections}